% Part: history
% Chapter: biographies
% Section: ernst-zermelo

\documentclass[../../../include/open-logic-section]{subfiles}

\begin{document}

\olfileid[hi]{his}{bio}{zer}

\olsection{अर्न्स्ट ज़र्मेलो}

\olphoto{zermelo-ernst}{अर्न्स्ट ज़र्मेलो}

अर्न्स्ट ज़र्मेलो का जन्म 27 जुलाई 1871 को बर्लिन, जर्मनी में हुआ। उनकी
पाँच बहनें थीं, किंतु परिवार खराब स्वास्थ्य से जूझता रहा और केवल तीन बहनें
वयस्कता तक जीवित रहीं। उनके बचपन में ही माता-पिता भी चल बसे, जिससे सत्रह
वर्ष की आयु में वे और उनके भाई-बहन अनाथ हो गए। ज़र्मेलो की कलाओं, विशेषतः
कविता, में गहरी रुचि थी। वे तीक्ष्ण, हाज़िरजवाब और आलोचनात्मक माने जाते थे।
उनकी सर्वाधिक प्रसिद्ध गणितीय उपलब्धियों में चयन अभिगृहीत की प्रस्तुति
(1904 में) और समुच्चय सिद्धांत का उनका अभिगृहीतीकरण (1908 में) हैं।

विश्वविद्यालय में ज़र्मेलो की रुचियाँ विविध थीं। उन्होंने भौतिकी, गणित और
दर्शन के पाठ्यक्रम लिए। हरमन श्वार्त्स के निर्देशन में ज़र्मेलो ने 1894 में
बर्लिन विश्वविद्यालय में अपना शोधप्रबंध \emph{Investigations in the
Calculus of Variations} पूरा किया। 1897 में उन्होंने गॉटिंगेन विश्वविद्यालय
में आगे अध्ययन करने का निर्णय लिया, जहाँ डेविड हिल्बर्ट के आधारभूत कार्य ने
उन्हें बहुत प्रभावित किया। 1899 में वे प्राध्यापक बनने के पात्र हो गए, किंतु
ग्यारह वर्ष बाद तक उन्हें पद नहीं मिला---संभवतः उनके विचित्र व्यवहार और
``व्यग्र उतावलेपन'' के कारण।

अंततः 1910 में ज़र्मेलो को ज़्यूरिख विश्वविद्यालय में सवेतन प्राध्यापक पद
मिला, किंतु तपेदिक के कारण 1916 में उन्हें सेवानिवृत्त होना पड़ा। स्वस्थ होने
के बाद 1921 में उन्हें फ़्राइबुर्ग विश्वविद्यालय में मानद प्राध्यापक पद दिया
गया। इस दौरान उन्होंने आधारभूत गणित पर काम किया। थोराल्फ़ स्कोलेम और कुर्ट
ग्योडेल की रचनाओं से वे चिढ़ गए और अपने शोधपत्रों में उनके दृष्टिकोणों की
सार्वजनिक आलोचना की। अलोकप्रियता और जर्मनी में हिटलर के सत्तारोहण के विरोध
के कारण उन्हें 1935 में फ़्राइबुर्ग के पद से हटा दिया गया।
 
ज़र्मेलो के जीवन के अंतिम वर्ष एकांत में बीते। 1935 में पद से हटाए जाने के
बाद उन्होंने गणित छोड़ दिया। वे देहात चले गए, जहाँ सादगी से रहे। 1944 में
उन्होंने विवाह किया और दृष्टि जाती रहने के कारण अपनी पत्नी पर पूर्णतः निर्भर
हो गए। 1951 तक ज़र्मेलो की दृष्टि पूरी तरह चली गई। उनका निधन 21 मई 1953 को
जर्मनी के ग्युंटरस्टाल में हुआ।

\begin{reading}
ज़र्मेलो की पूर्ण जीवनी के लिए \citet{Ebbinghaus2015} देखें। 1904 और 1908
के ज़र्मेलो के मूलभूत शोधपत्र मूल जर्मन में पढ़ने के लिए उपलब्ध हैं
\citep{Zermelo1904,Zermelo1908}। भौतिकी पर उनके लेखन सहित ज़र्मेलो की
संकलित रचनाएँ \citep{Ebbinghaus2010,Ebbinghaus2013} में अंग्रेज़ी अनुवाद
में उपलब्ध हैं।
\end{reading}

\end{document}
